\section{Ablations and failure modes}\label{sec:ablations}
Ablations use \NNabl{} further seeds (200--209), with the final configuration as reference on the same seeds
(Appendix Table~\ref{tab:ablations}). The intervals below are 95\% paired $t$-intervals of the difference from the
reference; with $n=10$, the smallest attainable exact $p$ is $0.002$.

\paragraph{Replication.}
On these seeds, \textsc{+Comp.} again has higher SO(3) closure ($+\NAbpSOthreeCompreferenceClosureMeandiff$,
$[\NAbpSOthreeCompreferenceClosureCilo,\NAbpSOthreeCompreferenceClosureCihi]$) and higher ground-truth distance
($+\NAbpSOthreeCompreferenceGtdistanceMeandiff$,
$[\NAbpSOthreeCompreferenceGtdistanceCilo,\NAbpSOthreeCompreferenceGtdistanceCihi]$) than BracketNet. Correct-closed
counts are \NAbSOthreeCompreferenceCc{} vs.\ \NAbSOthreeReferenceCc.

\paragraph{The anti-collapse term is essential; the schedule is not demonstrably needed.}
Lowering the covariance weight from 16 to 1 destroys recovery: the $T^2$ ground-truth distance rises by
$\NAbpTtwoWeakanticollapseGtdistanceMeandiff$, and correct-closed runs fall to \NAbTtwoWeakanticollapseCc{}. The
latent distorts (CKA with the true state \NAbTtwoWeakanticollapseCkatruth{} vs.\ \NAbTtwoReferenceCkatruth), while the
generators remain closed.

Replacing fit--close--refit by a constant penalty, or removing the generator freeze, changes neither SO(3)
ground-truth distance nor transport detectably:
\begin{itemize}
\item constant penalty: $[\NAbpSOthreeConstantpenaltyGtdistanceCilo,\NAbpSOthreeConstantpenaltyGtdistanceCihi]$ and
$[\NAbpSOthreeConstantpenaltyMseonezerostepCilo,\NAbpSOthreeConstantpenaltyMseonezerostepCihi]$;
\item no freeze: $[\NAbpSOthreeNofreezeGtdistanceCilo,\NAbpSOthreeNofreezeGtdistanceCihi]$ for ground-truth distance.
\end{itemize}
We therefore do not claim that staging is necessary.

Neither the Gram penalty nor the span-only closure objective makes a detectable difference. A three-times larger
$\lambda$ lowers closure slightly but leaves ground-truth distance and transport unchanged. A weaker $\lambda$
($0.3\times$) leaves more SO(3) runs non-closed.

\paragraph{Misspecified generator count.}
With $K{+}1$ generators no run recovers the true algebra in either group. For SO(3), closure regularization drives
incorrect closed spans in most runs (\NAbSOthreeKplusoneCc{} correct). This is forced. A closed 4-dimensional
subalgebra of $\so(4)$ is of type $\su(2)\oplus\mathfrak u(1)$, and the centralizer of the diagonal $\so(3)$ in
$\so(4)$ is trivial, so \emph{no} closed 4-dimensional span contains the true algebra. The learned spans capture on
average \NAbSOthreeKplusoneGtcoverage{} of it. With $K{-}1$, the learned algebra cannot contain the truth, and transport degrades sharply. Closure therefore
presupposes the right $K$; it cannot discover it.

\paragraph{Failure modes, summarized.}
\begin{enumerate}
\item \emph{Wrong closed class}: chiral $\su(2)$ for SO(3) (\NSOthreeBNChiral/20 fresh runs). Transport cannot
detect it: chiral runs fit transitions as well as correct ones.
\item \emph{Seed-level fitting failures} shared by all methods (Sec.~\ref{sec:limits}).
\item \emph{Latent distortion} without the anti-collapse term.
\item \emph{Misspecified $K$}.
\end{enumerate}
Only (1) is introduced by the closure objective itself. Proposition~\ref{prop:classify} predicts it, and the
participation-ratio diagnostic detects it without labels.
